%%%%%%%%%%%%%%%%%%%%%%%%%%%%%%%%%%%%%%%%%
% Medium Length Professional CV (Trey Hunner, LaTeXTemplates.com)
% Requires resume.cls in the same directory.
% Tailored for: Applied Scientist / Research Engineer Intern, Mistral AI
%%%%%%%%%%%%%%%%%%%%%%%%%%%%%%%%%%%%%%%%%

\documentclass{resume}

\usepackage[left=0.4in,top=0.4in,right=0.4in,bottom=0.4in]{geometry}
\usepackage[T1]{fontenc}
\usepackage[utf8]{inputenc}
\usepackage{charter}
\usepackage[hidelinks]{hyperref}

\name{Osinachi Chukwujama}
\address{Saint-Étienne, France}
\address{\href{mailto:osinachi.chukwujama@gmail.com}{osinachi.chukwujama@gmail.com} \\ +33 6 51 41 94 14}
\address{\href{https://github.com/vicradon}{github.com/vicradon} \\ \href{https://www.linkedin.com/in/chukwujama-osinachi/}{linkedin.com/in/chukwujama-osinachi}}

\begin{document}

%----------------------------------------------------------------------------------------
\begin{rSection}{Profile}
Photonics master's student and software engineer with a First Class engineering degree and several
years of production experience on large codebases. I implement generative models from first
principles in PyTorch, from the mathematical derivation through training and sampling, and I have
shipped models to edge GPUs and multimodal LLM features to production. Comfortable working
independently, writing fault-tolerant Python, and explaining technical work to any audience.
\end{rSection}

%----------------------------------------------------------------------------------------
\begin{rSection}{Research Interests}
Diffusion models, flow matching, score-based models, world models, optimisation.
\end{rSection}

%----------------------------------------------------------------------------------------
\begin{rSection}{Education}

{\bf Erasmus Mundus Joint Master, Intelligent Photonics (iPSRS)} \hfill {September 2026 -- 2028}\\
M1, Optics, Image, Vision and Multimedia, Université Jean Monnet, Saint-Étienne, France\\
\textit{Relevant coursework:} Fourier optics, digital image processing, probability and signal analysis, scientific computing with Python.

{\bf Bachelor of Engineering, Mechanical Engineering} \hfill {November 2017 -- August 2023}\\
Federal University of Technology, Owerri (FUTO), Nigeria. CGPA: 4.56/5.00 (First Class),\\
96th percentile in the department; best graduating student in Industrial and Production Engineering.

\end{rSection}

%----------------------------------------------------------------------------------------
\begin{rSection}{Skills}

\begin{tabular}{ @{} l @{\hspace{4ex}} p{0.72\textwidth} }
\textbf{Machine learning} & PyTorch: DDPM, OT-conditional flow matching, score matching, CNN/ResNet training, transfer learning (YOLOv5, ResNet18); ONNX export, TensorRT optimisation \\
\textbf{Mathematics} & Probability, variational inference (ELBO), continuity equation and optimal transport for flow matching, Fourier analysis \\
\textbf{Software} & Python (proficient), Go, TypeScript/JavaScript, PHP; Celery, RabbitMQ, PostgreSQL, GraphQL, Linux \\
\textbf{Engineering practice} & CI/CD pipelines, automated testing, typed codebases, logging and monitoring, backups and incident response \\
\textbf{Languages} & English: C1 (IELTS 8.0) \\
\end{tabular}

\end{rSection}

%----------------------------------------------------------------------------------------
\begin{rSection}{Research and ML Projects}

\begin{rSubsection}{Generative Modelling from First Principles}{2026}{Self-directed research}{}
\item Derived and implemented DDPM (forward/reverse process, ELBO, reparameterisation) and OT-conditional flow matching in PyTorch, covering training and sampling end to end.
\item Studied the links between score matching (denoising and implicit variants, NCSN), flow matching and normalising flows; wrote a forward-diffusion noising tool in Go for synthetic data generation.
\item Proposed a comparison of OT-CFM and DDPM on photonics/laser data (TPU Research Cloud application).
\end{rSubsection}

\begin{rSubsection}{Margin: Handwritten Notes with Multimodal AI Feedback}{2026}{Personal project, \href{https://margin.osinachi.me}{margin.osinachi.me}}{}
\item Built and deployed a pen-based note-taking app in which a multimodal LLM checks handwritten reasoning from images and returns immediate correctness feedback.
\item Switched LLM providers to cut feedback latency; notes sync to a Git repository through a pull-then-push workflow.
\end{rSubsection}

\begin{rSubsection}{Surveillance Drone with Real-Time Object Detection}{August 2023 -- March 2024}{Undergraduate final project and follow-up, supervised by Prof.\ O.O.\ Obiukwu}{MEE, FUTO}
\item Collected and annotated a three-class dataset, fine-tuned YOLOv5 (70\% mAP) and ran live inference on a drone's camera stream via ONNX.
\item Deployed the model on an NVIDIA Jetson Nano with TensorRT, cutting processing time by 44\%.
\end{rSubsection}

\end{rSection}

%----------------------------------------------------------------------------------------
\newpage
\begin{rSection}{Work Experience}

\begin{rSubsection}{Dstrct Group}{March 2025 -- May 2026}{Backend Engineer / Technical Lead}{Lagos, Nigeria}
\item Core engineer and technical lead on an enterprise intranet: a typed GraphQL backend (NestJS, TypeORM, PostgreSQL) of nearly 1,000 files spanning auth, documents, ticketing, chat, tasks, scheduling and WebSocket real-time features.
\item Worked directly with clients and stakeholders to set priorities and turn business requirements into engineering tasks.
\item Built an automated DOCX templating and PDF conversion pipeline generating formal documents from database records.
\end{rSubsection}

\begin{rSubsection}{Hotels.ng}{June 2024 -- January 2026}{Software/DevOps Engineer (part-time)}{Lagos, Nigeria}
\item Developed a highly concurrent Python/Celery application distributing AI workloads across workers over RabbitMQ.
\item Wrote CI/CD workflows for deployment, database migrations and automated testing.
\item Wrote logging and log-processing tooling that improved monitoring coverage by over 70\% versus manual checks.
\end{rSubsection}

\begin{rSubsection}{Generative CAD}{December 2021 -- June 2022}{Robotics \& 3D Printing Industrial Trainee}{Lagos, Nigeria}
\item Developed ROS path planning for robots, simulated in Gazebo and RViz; designed and 3D printed prototypes from customer requirements.
\end{rSubsection}

\end{rSection}

%----------------------------------------------------------------------------------------
\begin{rSection}{Publications}
\item \textit{Quantum Computing in Artificial Intelligence: A Review of Quantum Machine Learning Algorithms}, Path of Science.\\ \href{http://dx.doi.org/10.22178/pos.117-25}{doi:10.22178/pos.117-25}
\item \textit{Secure and Resilient Industrial IoT Architectures for Smart Manufacturing: A Comprehensive Review}, JERR (2025).\\ \href{https://doi.org/10.9734/jerr/2025/v27i61548}{doi:10.9734/jerr/2025/v27i61548}
\item \textit{Clinical Applications of 3D-Printed Medical Implants}, JOPIR.\\ \href{https://doi.org/10.69613/22dn3m33}{doi:10.69613/22dn3m33}
\end{rSection}

%----------------------------------------------------------------------------------------
\begin{rSection}{Teaching and Communication}
\item \textbf{Microsoft Certified Trainer} (Azure and AI): delivered enterprise workshops, including for MTN Nigeria, and prepared AI course materials.
\item \textbf{Former Microsoft Learn Student Ambassador}: trained over 300 students in software engineering and cloud technologies.
\item \textbf{Director of Research, Mechanical Engineering, FUTO}: ran hands-on training in machine learning, microcontrollers and 3D printing for 50+ students and staff. Co-founder of Hack4FUTO.
\item Technical writer at \href{https://vicradon.hashnode.dev/}{vicradon.hashnode.dev}.
\end{rSection}

\end{document}
